\begin{table}[htbp]
\centering
\caption{Cross-state entry-regime comparison: county-level sample, New Mexico and Colorado}
\label{tab:sumstats_state}
\footnotesize
\begin{tabular}{lcc}
\toprule
 & New Mexico & Colorado \\
\midrule
Entry regime & Free entry & Local option / capped \\
First recreational retail sales & 1 April 2022 & 1 January 2014 \\
\addlinespace
Counties (total) & 33 & 64 \\
\quad Border counties (geography) & 11 & 23 \\
\addlinespace
Counties in estimation sample & 32 & 42 \\
County--quarters & 360 & 530 \\
\quad Border county--quarters & 118 & 149 \\
\addlinespace
\multicolumn{3}{l}{\textit{County-quarter means}} \\
\quad Revenue per capita (\$) & 66.5 & 153.1 \\
\quad Dispensaries per 1,000 residents & 0.432 & 0.438 \\
\quad Revenue per dispensary (\$000s) & 213.2 & 524.8 \\
\quad County population & 65,955 & 125,571 \\
\quad Nonlegal pop.\ within 75 mi, border counties (000s) & 194.5 & 37.7 \\
\bottomrule
\end{tabular}
\begin{minipage}{0.86\textwidth}
\vspace{0.4em}\footnotesize
\textit{Notes:} County-quarter panel, 2022Q2 through 2025Q2, the sample used for the cross-state entry-regime comparison. Colorado's own market opened in 2014, but it is restricted to 2022Q2 onward.New Mexico is coded free entry; Colorado is coded locally capped because it delegates licensing to local governments and, as of 2019, 48.4 percent of its counties banned recreational cannabis businesses outright \citep{PayanBrownSong2021}.
\end{minipage}
\end{table}
